\section{Job Design and the Firm's Long Run Problem}
\label{app:jobdesign}

Section~\ref{sec:jobs_wages} studies the firm's AI deployment problem with the whole workflow assigned to a single worker whose wage is fixed, so that the firm chooses how to chain and automate steps in order to minimize the time production takes, $\sum_b \timecost{b}$.
This appendix develops the firm's long-run problem, in which it additionally bundles tasks into \emph{jobs}.
A single worker may perform several tasks, but must then hold all of their skills throughout, so the cost of a job is its total skill requirement times its total time.
Bundling economizes on the coordination costs incurred when work passes from one worker to another, at the expense of these broader skill requirements. 
The firm chooses the AI deployment and the job design jointly in the long-run problem to trade these forces off.
Figure~\ref{fig:hierarchical_production} of the main text illustrates the resulting three-layer structure of steps, tasks, and jobs in the long-run optimization problem.



\subsection{Jobs, Wages, and Hand-off Costs}
\label{sec:app.jobs}

This subsection restates, in full detail, the objects Section~\ref{sec:longrun.jobdesign} introduces, so that the long-run problem below is self-contained.

A job is a subsequence of tasks that are assigned to a single human worker.
Each task $T_b$ in the task sequence $\T = \{T_1, \dotsc, T_n\}$ has a time cost $\timecost{b}$, defined in Section~\ref{sec:model.tasks}, and a skill cost $\skillcost{b}$, introduced in Section~\ref{sec:longrun.jobdesign}; both depend on whether the task is completed manually or with the aid of AI.
A job $J$ is a contiguous subsequence of tasks, say $J = (T_b, \dotsc, T_{b+\ell})$ for some $b$ and $\ell$.
We say that a partition $\mathcal{J} = (J_1, \dotsc, J_p)$ of all tasks into jobs is a \emph{job design}.

The firm hires one worker for each job in its job design $\mathcal{J}$.
The worker for job $J_j$ completes all tasks associated with their job.
The total time needed to complete all tasks in job $J_j$ is $\sum_{T_b \in J_j} \timecost{b}$.

A worker's wage is determined by the total skill required to complete their job.
The total skill required to complete the tasks in job $J_j$ is $\sum_{T_b \in J_j} \skillcost{b}$.
We then assume that a worker's wage is proportional to the skill level of their job.
That is,
\begin{align}
\label{eq:wage_jobdesign}
\text{Wage}_j = \sum_{T_b \in J_j} \skillcost{b}.
\end{align}

The firm must pay workers their wage per unit of time regardless of which tasks they are assigned to complete at any given moment.
While we allow the firm to specialize its job design and partition its workflow into distinct jobs, we acknowledge that there are inherent inefficiencies and frictions that are introduced when splitting work among multiple workers.
Absent such frictions in our model, it would always be optimal for a firm to specialize its workers as much as possible, assigning a distinct worker to each task of its workflow. 
We therefore introduce \emph{hand-off costs} to a job design, as follows.
In addition to the time required completing tasks, a worker may need to spend time handing off their work to another worker if $J_j$ is not the final job in the job design $\mathcal{J}$.
We will assume throughout that the hand-off time of a worker assigned to job $J_j$ depends only on the final step of job $J_j$ (see red curly arrows in Figure \ref{fig:hierarchical_production}).
That is, conditional on where the hand-off occurs in the workflow, its cost does not otherwise depend on previous hand-off events.
Given step $s_i$, we write $\handofftime{i}$ to denote the additional hand-off time spent by a worker for whom the last step of their job is $s_i$.
For notational convenience we define the hand-off time for the final step $s_m$ to be zero: $\handofftime{m} = 0$.
Also for notational convenience, we will write $\handofftime{}(J_j)$ to denote the hand-off time of job $J_j$, which recall is equal to $\handofftime{i}$ if $s_i$ is the final step in job $J_j$.
Thus the total time needed to complete job $J_j$, including hand-off costs, is
\[
\handofftime{}(J_j) + \sum_{T_b \in J_j} \timecost{b}.
\]

Taking into account both time and skill requirements, the total wage bill paid to a worker assigned to job $J_j$, per unit of output, is
\begin{align}
\label{eq:wagebill}
\text{WageBill}_j = \Biggl(\sum_{T_b \in J_j} \skillcost{b} \Biggr) \Biggl(\handofftime{}(J_j) + \sum_{T_b \in J_j} \timecost{b} \Biggr).
\end{align}
Equation~\eqref{eq:wagebill} highlights two opposing forces that shape a job's boundaries.  
Adding more tasks to a worker's job increases the cumulative skill requirement and thus the wage rate for that worker.  
However, if tasks are kept separate, workers must incur additional hand-off costs at each job boundary.  
These two effects, higher wages from combining tasks versus higher hand-off costs from keeping them separate, jointly determine the optimal degree of task aggregation.  
Figure~\ref{fig:lr_job_design} in the main text illustrates these opposing forces in a minimal two-task workflow.
Section~\ref{sec:longrun.specialization} of the main text explores this trade-off in greater detail.

The notation for job design is collected in Panel~B of Table~\ref{tab:model_parameters} in the main text.


\subsection{Firm's Long-Run Optimization Problem}
\label{app:jobdesign.longrun}

In the long-run optimization problem, the firm anticipates the adjustment of worker wages to fit the skill requirements of each job and sets job responsibilities accordingly.
This contrasts with the AI deployment problem of Section~\ref{sec:jobs_wages}, which minimizes the time production takes while holding the assignment of steps to a single worker fixed. 
The long-run problem additionally optimizes how tasks are bundled into jobs.

Given the sequence of steps $\mathcal{S} = \{s_1, \dotsc, s_m\}$, the firm can design both the partition $\T$ of steps into tasks and the partition $\mathcal{J}$ of tasks into jobs.
The choice of $\T$ determines the time and skill requirements of each task.
Given $\T$, the choice of $\mathcal{J}$ then determines worker wage rates and time needed per unit of output (including hand-off costs).
Formally, we can write $\mathcal{P}(X)$ for the set of partitions of a sequence $X$ into contiguous subsequences.  Then the long-run optimization problem faced by the firm can be expressed as
\begin{align}
\label{eq:totalcost_with_handoff}
\min_{\T \in \mathcal{P}(\mathcal{S})} \
\min_{\mathcal{J} \in \mathcal{P}(\T)} \ 
\text{TotalCost}(\mathcal{J}; \T) 
= \sum_{J_j \in \mathcal{J}} \text{WageBill}_j = 
\sum_{J_j \in \mathcal{J}} \Biggl[\Biggl(\sum_{T_b \in J_j} \skillcost{b} \Biggr) \Biggl(\handofftime{}(J_j) + \sum_{T_b \in J_j} \timecost{b} \Biggr)\Biggr]
\end{align}
where recall that, for each $T_b \in \T$, $\timecost{b}$ is as described in Section~\ref{sec:model.tasks} and $\skillcost{b}$ as in Section~\ref{sec:longrun.jobdesign}, both depending on the selected mode of operation for individual steps.

This problem has a similar flavor to the fixed-boundary problem of Proposition~\ref{prop:mediumrun_dp}, with one additional layer: as steps are assigned to tasks, the resulting tasks must also be assigned to workers along the way by considering the hand-off costs incurred at the boundary of jobs.
Proposition~\ref{prop:totalcost_optimization_dp} of Section~\ref{sec:extensions.computation} states that, like that problem, the long-run problem admits a polynomial-time approximation scheme.
Equation~\eqref{eq:totalcost_with_handoff} is the objective \eqref{eq:lr_totalcost} that the proposition refers to, written out over the two partitions.
We prove the proposition here.

\begin{proof}
We prove the proposition following the same approach as the proof of Proposition~\ref{prop:mediumrun_dp}: we first state the cost minimization problem as a recursive problem and then use dynamic programming to approximate the optimal solution.
Since much of the construction carries over, we focus on what is new, the assignment of tasks to workers and the associated hand-off costs, and refer to the earlier proof for details that are unchanged.

\emph{Recursive formulation.}
Consider a workflow with $m$ steps $\mathcal{S} = \{s_1, s_2, \ldots, s_m\}$.
For $0 \leq i \leq m$, let $V(i)$ denote the minimum cost required to complete steps $1$ through $i$ using one or more jobs, including any hand-off costs to a subsequent worker.
Note then that $V(m)$ is the solution to the desired optimization problem, recalling that the hand-off time for the final job will always be $0$.
Note also that it is implicitly assumed in the definition of $V(i)$ that it optimizes over job designs for steps $1$ through $i$ in which the final job terminates after the completion of step $i$.

For job designs whose final job does not necessarily terminate after step $i$ but rather continues onward to subsequent steps, we extend the auxiliary function $R(i,\; \skillcost{},\; \timecost{})$ from the proof of Proposition~\ref{prop:mediumrun_dp}.
As before, $R(i,\; \skillcost{},\; \timecost{})$ denotes the minimum cost of completing steps $1$ through $i$ when the worker currently accumulating work, whom we call the ``active worker,'' has already been assigned tasks with total skill cost $\skillcost{}$ and total time cost $\timecost{}$, where $\timecost{}$ now also includes any hand-off cost for the active worker.
What is new is that the minimum is taken over job designs as well as AI deployment strategies: the recursion may either expand the active worker's job or close it and assign the remaining steps to other workers.
The base case $R(0,\; \skillcost{},\; \timecost{}) = \skillcost{}\,\timecost{}$ is unchanged, and as before we can think of step $i$ as the highest-indexed step that has not yet been assigned.

The two functions are linked, for all $0 \leq i \leq m$, by
\[ V(i) = R(i,\; 0,\; \handofftime{i}). \]
This is because, in $V(i)$, all job designs must end with a job that completes at step $i$ and hands off to the following worker.
Incurring this hand-off cost is equivalent to assigning the active worker an additional task with zero skill requirement and a time cost of $\handofftime{i}$ right after step $i$ itself, which is exactly the scenario $R(i,\; 0,\; \handofftime{i})$ represents.

For each $i \geq 1$, the minimum cost function $R(i,\; \skillcost{},\; \timecost{})$ satisfies the recursive relation
\begin{align*}
R(i,\; \skillcost{},\; \timecost{}) = \min\Biggl\{
& \ \underbrace{\skillcost{}\,\timecost{} + V(i)}_{\text{job closed}}, \ \underbrace{R\bigl(i-1,\; \skillcost{} + \manualSkill{i},\; \timecost{} + \manualTime{i}\bigr)}_{\text{Step $i$ is manual}}, \ \underbrace{\min_{0 \le r < i}\ R\!\left(r,\; \skillcost{} + \AIskill{i},\; \timecost{} + \frac{\AItime{i}}{\prod_{s=r+1}^{i} q_s}\right)}_{\text{Step $i$ optimally chained}}
\Biggr\}.
\end{align*}
In words, when assigning step $i$, the three options the firm considers are:
\begin{itemize}
    \item \textbf{Job closure.}
    The term $\skillcost{}\,\timecost{} \; + \; V(i)$ corresponds to not adding any further steps to the job of the active worker.
    All steps $1$ through $i$ will be completed by other workers.
    The active worker's total cost is then $\skillcost{}\,\timecost{}$, and $V(i)$ is the total cost of completing steps $1$ through $i$ (including hand-off costs to the active worker).
    This option is the margin new to the long run.

    \item \textbf{Manual step.}
    The term $R(i-1,\; \skillcost{} + \manualSkill{i},\; \timecost{} + \manualTime{i})$ corresponds to designating step $i$ for manual completion, and adding the corresponding (singleton) task.
    As in Appendix~\ref{app:mediumrun_dp}, the manual execution costs $(\manualSkill{i}, \manualTime{i})$ of step $i$ are added to the accumulated costs of tasks already assigned to the active worker, extending their job to include step $i$ as well.

    \item \textbf{AI chain.}
    The term $\min_{0 \le r < i} \ R\left(r,\; \skillcost{} + \AIskill{i},\; \timecost{} + \frac{\AItime{i}}{\prod_{s=r+1}^{i} q_s}\right)$ corresponds to creating an AI chain task for steps $r+1$ through $i$, with step $i$ being augmented and steps $r+1$ through $i-1$ automated, and assigning the resulting task to the job of the active worker.
    Step $i$ is thus chained with zero or more preceding (automated) steps, and the associated AI management costs are added to the active worker's skill and time costs.
    The cut $r$ is the highest-indexed step that is not added to this AI chain, which can be any value between $0$ and $i-1$.
\end{itemize}
Note that in evaluating the AI chain option, step $i$ can only be augmented (rather than automated), as each step of the recursion adds a full AI chain to the active worker's job; AI strategies in which step $i$ is automated are considered when performing the calculations $R(r,\; \skillcost{},\; \timecost{})$ with $r > i$.
The cost of the optimal joint AI strategy and job design for the full sequence of $m$ steps is then $V(m) = R(m,\; 0,\; 0)$.

\emph{Discretization and Computation Time.}
As in the proof of Proposition~\ref{prop:mediumrun_dp}, we fill in the table of values $R(i,\; \skillcost{},\; \timecost{})$ by dynamic programming after discretizing the skill and time coordinates to powers of $(1+\epsilon)$; only the range of the discretization requires a new argument, to account for hand-off costs.
Completing every step manually with a separate worker incurs, for each step, a time and hand-off cost of at most $B$ each and a skill cost of at most $B$, for a total cost of at most $2B^2$ per step, so no solution containing a job with total cost above $2mB^2$ need be considered.
Since skill costs combine additively and every non-zero skill or time cost is at least $1/B$, it suffices to tabulate skill costs in $[1/B,\, mB]$ and time costs in $[1/B,\, 2mB^3]$.
Each range contains $O(\epsilon^{-1}\log(mB))$ powers of $(1+\epsilon)$, so the table has $O(m\,\epsilon^{-2}\log^2(mB))$ entries, each filled in time $O(m)$ via the recursion above, for a total runtime of $O(m^2\,\epsilon^{-2}\log^2(mB))$; recording which option attains each minimum lets the firm read the optimal AI strategy and job design off $V(m) = R(m,\; 0,\; 0)$.
The discretization error is bounded exactly as before: skill and time costs accumulate additively along the manual and chaining options, so rounding down to powers of $(1+\epsilon)$ maintains a multiplicative error of at most $(1+\epsilon)$ in each coordinate, while each job's wage bill is formed by the single multiplication $\skillcost{}\,\timecost{}$ when the first option closes it, contributing a factor of at most $(1+\epsilon)^2$.
Rescaling $\epsilon$ by a constant therefore yields the claimed $(1+\epsilon)$ approximation.
\end{proof}

\subsection{Connection to the Division of Labor Literature}

The job-design problem developed above connects to the literature on the division of labor within a firm.
Applying the Coase--Williamson logic that firm boundaries are shaped by coordination, transaction, and governance costs \citep{coase1937nature, williamson1971vertical, williamson1979transaction} to the ``boundary of the job,'' the firm chooses how many and which tasks to bundle into a worker's role, determining the degree of worker specialization endogenously \citep{murphy1986specialization, becker1992division, dessein2006adaptive}.
The hand-off costs in our model are an intra-firm analogue of transaction costs that AI does not directly reduce, since they capture coordination frictions arising from tacit or social knowledge that is not easily executable by AI \citep{deming2017growing}.\footnote{
\cite{ide2025ai} provides a related view on how AI affects the structure of work, embedding AI as an autonomous problem solver in a knowledge hierarchy model \citep{garicano2000hierarchies, garicano2006organization} and studying how it reshapes spans of control within the firm.
}
Even so, AI chaining restructures the division of labor on two margins: on the intensive margin, by pulling previously manual steps ``under the hood'' of a single AI-executed task; and on the extensive margin, by lowering the cost of linking adjacent steps across a job boundary and thereby making it optimal to redraw that boundary \citep{autor2025expertise}.
